\chapter{因果生存分析}

% 生存曲线坐标框（本章四幅生存曲线图共用）
\newcommand{\survaxes}{%
  \draw (0,0) rectangle (9,5);
  \foreach \m in {0,12,...,120}{%
    \draw (\m*0.075,0) -- (\m*0.075,-0.10);
    \node[below,font=\scriptsize,inner sep=1.5pt] at (\m*0.075,-0.10) {\m};}
  \foreach \p in {0.5,0.6,0.7,0.8,0.9,1.0}{%
    \draw (0,{(\p-0.5)*10}) -- (-0.10,{(\p-0.5)*10});
    \node[left,font=\scriptsize,inner sep=1.5pt] at (-0.12,{(\p-0.5)*10}) {\p};}
  \node[below=16pt,font=\small] at (4.5,-0.10) {随访月数};
  \node[rotate=90,above=20pt,font=\small] at (-0.78,2.5) {生存概率};}

在此前的各章中，我们关心的是处理对某一特定时点发生的结局的因果效应。例如，我们估计过戒烟对 1982 年测得的体重增加的影响。然而，许多因果问题关心的是处理对所关注事件发生时间的影响。例如，我们可能想估计戒烟对死亡时间的影响——无论死亡何时发生。这是生存分析的一个例子。

使用“生存”一词，并不意味着所关注的事件必须是死亡。“生存分析”这一术语（与之等价的说法是“失效时间分析”）适用于一切关于事件发生时间的分析，其中的事件可以是死亡、结婚、入狱、癌症、流感感染等等。生存分析需要一些特殊的考量与技术，因为许多个体的失效时间可能发生在研究结束之后，因而无从知晓。本章在\emph{时不变治疗}的简化设定下，概述生存分析的基本技术。

\section{风险率与风险}

假设我们想估计戒烟 $A$（1：是，0：否）对死亡时间 $T$ 的平均因果效应，时间自随访开始起计。这是生存分析的一个例子：结局是所关注事件的发生时间，而该事件可在随访开始后的任何时刻发生。在多数随访研究中，并非全部、甚至并非多数个体的所关注事件都会被观测到发生。原因在于，多数随访研究都有一个日期，此后便再没有任何个体的信息：随访的行政终止。

随访的行政终止之后，无法再使用任何新增数据。在随访的行政终止之前未发生所关注事件的个体，其生存时间被行政删失，也就是说，我们知道他们活过了随访的行政终止，但不知道他们此后还活了多久。例如，设我们在 NHEFS 中 1629 名基线年龄 25–74 岁、且活过 1982 年的吸烟者中做上述生存分析。对全部个体，随访开始于 1983 年 1 月 1 日，随访的行政终止为 1992 年 12 月 31 日。我们把行政删失时间定义为随访行政终止日期与随访开始日期之差。在本例中，由于所有人的随访开始时间与随访行政终止时间都相同，这一时间对全部个体相同——120 个月。这 1629 人中，只有 318 人在 1992 年底之前死亡，因此其余 1311 人的生存时间被行政删失。

\mnote{在错时进入（即随访开始日期因人而异）的研究中，即便随访的行政终止日期对所有人相同，不同个体的行政删失时间也会不同。}

行政删失是生存分析固有的问题——关于戒烟与死亡的研究很少（如果有的话）会把一个队列随访到全部个体死亡——但行政删失并不是生存分析中可能出现的唯一一类删失。与任何其他因果分析一样，生存分析可能还需要处理非行政类型的删失，例如失访（如退出研究）与竞争事件（见精要 17.1）。前几章讨论过如何通过标准化或逆概率加权来校正非行政删失引入的选择偏倚；同样的方法可用于生存分析。因此，本章集中讨论行政删失。对非行政删失更细致的考察留到本书第三部分，因为非行政删失一般是一个时变过程，而行政删失的时间在基线时就已固定。

\begin{finepoint}[竞争事件]
如 8.5 节所述，\emph{竞争事件}是指使所关注事件（如中风）无法发生的事件（通常是死亡）：死于其他原因（比如癌症）的人永远不可能发生中风。在生存分析中，关键决策在于是否把竞争事件视为一种非行政删失。
\begin{itemize}
\item 若把竞争事件视为删失事件，则分析实际上是在试图模拟这样一个总体：其中死于其他原因要么被消除，要么被弄得与中风的风险因子无关。所得效应估计难以解释，且可能并不对应一个有意义的估计目标（见第 8 章）。此外，这种删失可能在零假设下引入选择偏倚，从而需要用所关注事件之已测风险因子的数据做校正（比如通过逆概率加权）。
\item 若不把竞争事件视为删失事件，则分析实际上把事件发生时间设为无穷大。也就是说，死者被认为在其死亡到随访行政终止之间发生中风的概率为零。处理对中风之效应的估计难以解释，因为非零的估计可能来自处理对死亡的直接效应，而死亡会阻止中风的发生。
\end{itemize}
处理竞争事件的另一种办法，是构造一个同时包含竞争事件与所关注事件的复合事件（如死亡与中风），并对该复合事件做生存分析。这一做法实际上消除了竞争事件，但从根本上改变了因果问题。同样，所得效应估计难以解释，因为非零的估计可能来自处理对中风或对死亡的效应。还有一种办法，是把推断限制在无论接受何种治疗水平都不会死亡的那部分个体的主分层上。这一做法瞄准的是一种局部平均效应，如第 16 章所定义，它使解释与有效估计都格外困难。

上述策略都没有为竞争事件问题提供令人满意的解决方案。确实，竞争事件的存在会对因果估计目标的含义提出逻辑上的质疑，这是统计技术无法回避的。关于处理竞争事件的种种方法及其挑战的详细描述，见 Young 等（2019）的讨论。最近，Stensrud 等（2020，2021）提出了一种基于可分离效应的方法，该概念见技术注记 23.3。
\end{finepoint}

在本例中，死亡月份 $T$ 的取值可从 1（1983 年 1 月）到 120（1992 年 12 月）。对随访期间死亡的 102 名受治疗者（$A=1$）与 216 名未受治疗者（$A=0$），$T$ 是已知的；对其余 1311 人，$T$ 被行政删失（即我们只知道它大于 120 个月）。因此，我们无法像前几章处理所关心的结局那样计算平均生存时间 $\E[T]$。相反，在生存分析中，我们需要使用其他能够容纳行政删失的度量。

\mnote{为简单起见，我们假定凡未经确认死亡者都活过了随访期。现实中，有些人可能已经死亡，但无法确认（例如缺少死亡证明或代答访谈）。同样为简单起见，我们忽略精要 12.1 描述的问题。}

一些常用的度量是生存概率、风险与风险率。下面定义这些量，它们都是生存时间 $T$ 的函数。

生存概率 $\Pr[T>k]$，或简称第 $k$ 月的生存，是活过时间 $k$ 的个体所占的比例。若我们计算从第 1 月直到随访行政终止 $k_{\text{end}}=120$ 的各月生存，并把它们画在一条水平时间轴上，就得到\emph{生存曲线}。生存曲线从 $k=0$ 时的 $\Pr[T>0]=1$ 出发，随后随 $k=1,2,\dots,k_{\text{end}}$ 单调下降——即不上升。或者，我们可以把时间 $k$ 处的\emph{风险}，或称\emph{累积发病}，定义为 1 减去生存：$1-\Pr[T>k]=\Pr[T\le k]$。累积发病曲线从 $\Pr[T\le 0]=0$ 出发，在随访期内单调上升。

在生存分析中，量化处理效应的一种自然做法，是比较各治疗水平下、某个或某些时点 $k$ 的生存或风险。\mnote{可由生存曲线导出的其他效应度量，是寿命损失年与限制平均生存时间。}当然，在我们戒烟的例子中，这种曲线对比未必具有因果解释，因为受治疗者与未受治疗者很可能不可交换。不过，暂且（实际上直到 17.4 节）假设戒烟者（$A=1$）与未戒烟者（$A=0$）边缘可交换。那么我们就可以构造图 17.1 所示的生存曲线，并在所有时点 $k$ 上比较 $\Pr[T>k\given A=1]$ 与 $\Pr[T>k\given A=0]$。例如，120 个月时的生存，在戒烟者中为 76.2\%，在未戒烟者中为 82.0\%。或者，我们也可以对比风险而非生存。例如，120 个月时的风险，在戒烟者中为 23.8\%，在未戒烟者中为 18.0\%。\mnote{图 17.1}

\begin{figure}[htbp]
\centering
\begin{tikzpicture}[x=1cm,y=1cm,causal]
  \survaxes
  \draw[thick] (0.00,5.00) -- (0.45,4.92) -- (0.90,4.85) -- (1.35,4.76) -- (1.80,4.68) -- (2.25,4.60) -- (2.70,4.50) -- (3.15,4.39) -- (3.60,4.28) -- (4.05,4.18) -- (4.50,4.09) -- (4.95,4.01) -- (5.40,3.93) -- (5.85,3.82) -- (6.30,3.70) -- (6.75,3.63) -- (7.20,3.56) -- (7.65,3.47) -- (8.10,3.38) -- (8.55,3.29) -- (9.00,3.20);
  \draw[thick,dashed] (0.00,5.00) -- (0.45,4.88) -- (0.90,4.78) -- (1.35,4.66) -- (1.80,4.55) -- (2.25,4.42) -- (2.70,4.30) -- (3.15,4.14) -- (3.60,4.00) -- (4.05,3.88) -- (4.50,3.78) -- (4.95,3.68) -- (5.40,3.58) -- (5.85,3.40) -- (6.30,3.22) -- (6.75,3.08) -- (7.20,2.95) -- (7.65,2.84) -- (8.10,2.75) -- (8.55,2.68) -- (9.00,2.62);
  \node[font=\small] at (7.35,3.85) {$A=0$};
  \node[font=\small] at (7.35,2.52) {$A=1$};
\end{tikzpicture}
\caption{}
\label{fig:17.1}
\end{figure}

在任何时点 $k$，我们还可以计算在 $k$ 之前尚未发生该事件者中、发生该事件的比例。这就是\emph{风险率} $\Pr[T=k\given T>k-1]$。严格说来，这是离散时间风险率，即时间以离散区间度量（而非连续度量）的研究中的风险率。由于在现实研究中，时间确实以离散区间（年、月、周、日……）而非真正连续的方式度量，下文我们把离散时间风险率简称为风险率。

风险与风险率是不同的度量。风险的分母——基线时的个体数——在各时点 $k$ 上保持不变，其分子——从基线到 $k$ 之间的全部事件——是累积的。也就是说，随着 $k$ 增大，风险要么持平、要么上升。相反，风险率的分母——在 $k$ 时仍存活的个体数——随时间变化，其分子只包含近期事件，即区间 $k$ 内的事件。也就是说，风险率可随时间上升或下降。在本例中，120 个月时的风险率在戒烟者中为 0\%（因为该组最后一次死亡发生在 113 个月），在未戒烟者中为 $1/986=0.10\%$；并且 0 到 120 个月之间的风险率曲线大致呈字母 M 的形状。\mnote{代码：程序 17.1。用于比较生存曲线的常规统计检验（对数秩检验）给出的 $P$ 值为 0.005。}

在生存分析中，量化处理效应的一种常见做法，是估计受治疗者与未受治疗者风险率之比，即\emph{风险率比}。然而，如精要 17.2 所述，风险率比存在一些问题。因此，本书的生存分析更看重生存／风险，而非风险率。但这并不意味着应当完全忽略风险率：估计风险率往往是估计生存与风险的有用中间步骤。

\section{由风险率到风险}

在生存分析中，整理分析数据集主要有两种方式。第一种方式下，数据库的每一行对应一个人。这种数据格式——在存在时变治疗与混杂因子时常被称为“宽”格式——是本书迄今使用的格式。上一节的分析中，数据集有 1629 行，每人一行。

第二种方式下，数据库的每一行对应一个人时。也就是说，第一行包含 1 号个体在 $k=0$ 时的信息，第二行包含 1 号个体在 $k=1$ 时的信息，第三行包含 1 号个体在 $k=2$ 时的信息，依次类推，直到 1 号个体的随访结束。下一行则包含 2 号个体在 $k=0$ 时的信息，等等。这种人时（或“长”）数据格式，是本章多数生存分析、以及第三部分一切含时变治疗的分析所使用的格式。在我们的戒烟例子中，人时数据集有 176\,764 行，每人月一行。

在人时数据格式中，为编码截至 $k$ 的生存信息，定义一个时变的事件指示变量 $D_k$ 是有用的。对每个人、每个月 $k$，若 $T\le k$ 则 $D_k$ 取 1，若 $T>k$ 则取 0。图 17.2 的因果图显示了处理 $A$ 以及时点 1 与时点 2 的事件指示变量 $D_1$ 与 $D_2$。变量 $U$ 表示（一般未测量的）增加所关注事件易感性的因素。注意，这些易感因素有时也是时变的；此时它们可在因果图中画成 $U_0,U_1,\dots$ 等等。第三部分处理治疗本身为时变的情形。

\begin{figure}[htbp]
\centering
\begin{tikzpicture}[causal,every node/.style={font=\small}]
  \node (A)  at (0,0)    {$A$};
  \node (D1) at (2.2,0)  {$D_1$};
  \node (D2) at (4.4,0)  {$D_2$};
  \node (U)  at (-1.4,-1.7) {$U$};
  \draw[->] (A) -- (D1);
  \draw[->] (D1) -- (D2);
  \draw[->] (A) to[out=55,in=125] (D2);
  \draw[->] (U) -- (D1);
  \draw[->] (U) -- (D2);
\end{tikzpicture}
\caption{}
\label{fig:17.2}
\end{figure}

在人时数据格式中，某个体在时间 $k$ 的那一行包含指示变量 $D_{k+1}$。\mnote{图 17.2}在本例中，人时数据集的第一行是 1 号个体在 $k=0$ 时的信息，包含 $D_1$：若该个体在第 1 月内死亡则为 1，否则为 0；第二行是 1 号个体在 $k=1$ 时的信息，包含 $D_2$：若该个体在第 2 月内死亡则为 1，否则为 0；依此类推。每个个体在数据集中的最后一行，要么是其第一个满足 $D_{k+1}=1$ 的行，要么是对应第 119 月的行。

\mnote{按定义，人人都必须活过第 0 月才能进入数据集，即对全部个体都有 $D_0=0$。}

利用时变结局变量 $D_k$，可以把 $k$ 时的生存定义为 $\Pr[D_k=0]$（等于 $\Pr[T>k]$），把 $k$ 时的风险定义为 $\Pr[D_k=1]$（等于 $\Pr[T\le k]$）。$k$ 时的风险率定义为 $\Pr[D_k=1\given D_{k-1}=0]$。对 $k=1$，风险率等于风险，因为按定义，所有人在 $k=0$ 时都存活。

$k$ 时的生存概率，是活过 0 到 $k$ 之间每一区间的条件概率之积。例如，$k=2$ 时的生存 $\Pr[D_2=0]$，等于 $k=1$ 时的生存概率 $\Pr[D_1=0]$ 乘以在活过 $k=1$ 的条件下 $k=2$ 时的生存概率 $\Pr[D_2=0\given D_1=0]$。更一般地，$k$ 时的生存为
\[
\Pr[D_k=0]=\prod_{m=1}^{k}\Pr[D_m=0\given D_{m-1}=0]
\]
也就是说，$k$ 时的生存，等于此前所有时点“1 减风险率”之积。若已知截至 $k$ 的风险率，就能轻易算出 $k$ 时的生存（或 $k$ 时的风险，它不过是 1 减生存）。

$k$ 时的风险率 $\Pr[D_k=1\given D_{k-1}=0]$ 可用非参数方法估计：用区间 $k$ 内的病例数除以区间 $k-1$ 结束时的存活人数。若把这个估计代入上述公式，所得的 $k$ 时生存 $\Pr[D_k=0]$ 的非参数估计，就称为 Kaplan-Meier 估计量，或称乘积极限估计量。图 17.1 就是用 Kaplan-Meier 估计量构造的；只要随访期内的失效总数足够大，它就是生存曲线的优良估计量。典型情形下，各区间内的病例数很少（甚至为零），因而 $k$ 时风险率 $\Pr[D_k=1\given D_{k-1}=0]$ 的非参数估计会非常不稳定。若我们关心的是某个特定 $k$ 处风险率的估计，可能需要通过参数模型进行平滑（见第 11 章与精要 17.3）。

\begin{finepoint}[风险率比的隐患]
把风险率比用作因果效应的度量时，需要注意风险率比的两个重要性质。

第一，由于风险率随时间变化，风险率比一般也随时间变化。也就是说，$k$ 时的比值可能与 $k+1$ 时的不同。然而，许多已发表的生存分析只报告一个风险率比，这通常是拟合 Cox 比例风险模型的结果——该模型忽略与时间的交互作用，假定风险率比恒定。所报告的风险率比是各 $k$ 特定风险率比的加权平均，这使其难以解释。若风险罕见、且删失只发生在一个共同的行政删失时间 $k_{\text{end}}$，则 $k$ 时风险率比的权重，与未受治疗者中在 $k$ 时发生的全部事件数成正比。（严格说来，权重等于给定 $A=0$ 与 $T<k_{\text{end}}$ 时 $T$ 在 $k$ 处的条件密度。）由于是加权平均，即便生存曲线并不相同，所报告的风险率比也可能为 1。与“那个”风险率比不同，生存概率与风险的比和差都是相对于固定时间区间定义的，例如 5 年生存差、120 个月风险比。

第二，即便给出各时点特定的风险率比，其因果解释也并不直接。设治疗在 $k$ 时之前杀死了所有高危个体，而对其他人没有效应。那么 $k+1$ 时的风险率比，比较的是活过 $k$ 的受治疗者与未受治疗者。在治疗组中，存活者全是低危个体（因为高危者已被治疗杀死）；在未治疗组中，存活者是高危与低危个体的混合（因为治疗没有剔除前者）。结果，即便治疗对任何个体都无益，$k+1$ 时的风险率比也会小于 1。

这一表面上的悖论，是以受治疗影响的事后变量（即在 $k$ 时存活）为条件所致选择偏倚的一个例子。例如，2 时的风险率比，是活过 1 时者中 2 时发生事件的条件概率 $\Pr[D_2=1\given D_1=0,A]$。如图 17.3 的因果图所示，以对撞节点 $D_1$ 为条件，一般会打开路径 $A\to D_1\leftarrow U\to D_2$，从而在 $D_1=0$ 者中诱导出治疗 $A$ 与事件 $D_2$ 之间的关联。若受治疗者与未受治疗者的生存曲线相同，即若没有从 $A$ 指向事件指示变量的箭头，则风险率比这种内在的选择偏倚不会发生。Hernán（2010）描述过这一问题的一个例子。
\end{finepoint}

\begin{figure}[htbp]
\centering
\begin{tikzpicture}[causal,every node/.style={font=\small}]
  \node (A)  at (0,0)    {$A$};
  \node[draw,inner sep=2.5pt] (D1) at (2.2,0) {$D_1$};
  \node (D2) at (4.4,0)  {$D_2$};
  \node (U)  at (-1.4,-1.7) {$U$};
  \draw[->] (A) -- (D1);
  \draw[->] (D1) -- (D2);
  \draw[->] (U) -- (D1);
  \draw[->] (U) -- (D2);
\end{tikzpicture}
\caption{}
\label{fig:17.3}
\end{figure}

对风险率做参数估计的一种简便方法，是对 $\Pr[D_{k+1}=1\given D_k=0]$ 拟合一个 logistic 回归模型，该模型在每个 $k$ 处仅限于活过 $k$ 的个体。\mnote{图 17.3}使用人时数据格式时，这一模型的拟合十分简便。\mnote{尽管每个人在人时数据中出现于多行，只要风险率模型正确，常规 logistic 回归程序输出的参数估计之标准误就是正确的。}在本例中，我们可以通过拟合如下 logistic 模型来估计受治疗者与未受治疗者的风险率
\[
\operatorname{logit}\Pr[D_{k+1}=1\given D_k=0,A]=\theta_{0,k}+\theta_1A+\theta_2A\times k+\theta_3A\times k^2
\]
其中 $\theta_{0,k}$ 是时变截距，可用时间的某个灵活函数来估计，例如 $\theta_{0,k}=\theta_0+\theta_4k+\theta_5k^2$。灵活的时变截距允许风险率随时间变化，处理 $A$ 与时间之间的乘积项（$\theta_2A\times k+\theta_3A\times k^2$）则允许风险率比随时间变化。关于 logistic 模型如何近似风险率模型，详见技术注记 17.1。除 logit 之外的函数（如 probit）也可用于对二值结局建模，因而也可用于估计风险率。

随后，我们分别对受治疗者与未受治疗者，把 logistic 模型给出的 $\Pr[D_{k+1}=1\given D_k=0,A=a]$ 之估计取“1 减”再相乘，即得生存 $\Pr[D_{k+1}=0\given A=a]$ 的估计。

\begin{finepoint}[生存分析模型]
生存分析的方法需要容纳因随访行政终止而预期发生的失效时间删失。

生存分析的非参数方法，如构造 Kaplan-Meier 曲线，对因行政删失而未观测到的失效时间之分布不作任何假定。相反，生存分析的参数模型则假定失效时间或风险率服从某种特定的统计分布（如指数分布、Weibull 分布）。正文中用于估计风险率的 logistic 模型就是参数模型的一例。

其他生存模型，如 Cox 比例风险模型与 17.6 节中的加速失效时间（AFT）模型，并不假定失效时间或风险率服从特定分布。这些模型对模型中全部协变量取零值时的风险率形状（常称为基线风险率）不作假定。然而，这些模型对基线风险率与其他协变量取值组合下的风险率之间的关系施加了先验限制。因此，这些方法被称为半参数方法。

应用生存分析的综述见 Hosmer、Lemeshow 与 May（2008）的书。更形式化的描述见 Fleming 与 Harrington（2005）、Kalbfleisch 与 Prentice（2002）的书。
\end{finepoint}

\begin{figure}[htbp]
\centering
\begin{tikzpicture}[x=1cm,y=1cm,causal]
  \survaxes
  \draw[thick] (0.00,5.00) -- (0.90,4.90) -- (1.80,4.72) -- (2.70,4.55) -- (3.60,4.36) -- (4.50,4.15) -- (5.40,3.97) -- (6.30,3.75) -- (7.20,3.52) -- (8.10,3.32) -- (9.00,3.18);
  \draw[thick,dashed] (0.00,5.00) -- (0.90,4.82) -- (1.80,4.60) -- (2.70,4.33) -- (3.60,4.05) -- (4.50,3.80) -- (5.40,3.58) -- (6.30,3.33) -- (7.20,3.10) -- (8.10,2.90) -- (9.00,2.68);
  \node[font=\small] at (7.35,3.85) {$A=0$};
  \node[font=\small] at (7.35,2.52) {$A=1$};
\end{tikzpicture}
\caption{}
\label{fig:17.4}
\end{figure}

图 17.4 显示了对风险率做参数估计后得到的生存曲线。\mnote{图 17.4}这些曲线是图 17.1 中曲线的平滑版本。

这一程序的有效性要求风险率模型无错误设定。在本例中，这一假定似乎是合理的，因为我们以完全非参数的方式估计生存时，得到的生存估计与上一节基本相同。生存估计的 95\% 置信区间，可通过对总体中的个体做自助法轻易构造。

\section{删失为何重要}

在我们的研究中，唯一的删失来源是一个对所有个体都相同的行政删失时间 $k_{\text{end}}=120$。在这种简单设定下，上一节所述的生存估计程序可谓杀鸡用牛刀。人们可以直接用接受治疗 $a$ 且存活到 $k+1$ 的个体所占比例来估计生存概率 $\Pr[D_{k+1}=0\given A=a]$，或者对每个时点 $k=0,1,\dots,k_{\text{end}}$ 分别拟合 $\Pr[D_{k+1}=0\given A]$ 的 logistic 模型。

现在设想个体在不同日期开始随访——即研究存在错时进入——但随访的行政终止日期对所有人相同。由于行政删失时间是随访行政终止与随访开始时间之差，不同个体的行政删失时间将不同。在这种设定下，定义一个截至 $k$ 的删失时变指示变量 $C_k$ 是有用的。对每个人、每个月 $k$，若随访的行政终止大于 $k$ 则 $C_k$ 取 0，否则取 1。在人时数据格式中，某个体在时间 $k$ 的那一行包含指示变量 $C_{k+1}$。在我们的数据集中没有纳入这个变量，因为在 120 个月之前的所有时点 $k$ 上，对所有个体都有 $C_{k+1}=0$。在行政删失随机（即因人而异）的一般情形下，指示变量 $C_{k+1}$ 会在不同时点 $k$ 对不同个体由 0 转为 1。

\begin{techpoint}[用 logistic 模型近似风险率比]
在风险率模型 $\Pr[D_{k+1}=1\given D_k=0,A]=\Pr[D_{k+1}=1\given D_k=0,A=0]\times\exp(\alpha_1A)$ 中，全部时点 $k+1$ 上的（离散时间）风险率比
\[
\frac{\Pr[D_{k+1}=1\given D_k=0,A=1]}{\Pr[D_{k+1}=1\given D_k=0,A=0]}
\]
为 $\exp(\alpha_1)$。对等式两边取对数，得 $\log\Pr[D_{k+1}=1\given D_k=0,A]=\alpha_{0,k}+\alpha_1A$，其中 $\alpha_{0,k}=\log\Pr[D_{k+1}=1\given D_k=0,A=0]$。

设 $k+1$ 时的风险率很小，即 $\Pr[D_{k+1}=1\given D_k=0,A]\approx0$。则 $k+1$ 时的“1 减风险率”接近 1，风险率近似等于优势：$\Pr[D_{k+1}=1\given D_k=0,A]\approx\Pr[D_{k+1}=1\given D_k=0,A]/\Pr[D_{k+1}=0\given D_k=0,A]$。于是有
\[
\log\frac{\Pr[D_{k+1}=1\given D_k=0,A]}{\Pr[D_{k+1}=0\given D_k=0,A]}=\operatorname{logit}\Pr[D_{k+1}=1\given D_k=0,A]\approx\alpha_{0,k}+\alpha_1A
\]
也就是说，若 $k+1$ 时的风险率接近零，我们可以在正文所用的那种 logistic 模型 $\operatorname{logit}\Pr[D_{k+1}=1\given D_k=0,A]=\theta_{0,k}+\theta_1A$ 中，用 $\theta_1$ 近似对数风险率比 $\alpha_1$（Thompson 1977）。经验上，当所有 $k$ 都有 $\Pr[D_{k+1}=1\given D_k=0,A]<0.1$ 时，这一近似通常被认为足够准确。

这一罕见事件条件几乎总能保证成立：我们只需把时间单位 $k$ 定义得足够短，使 $\Pr[D_{k+1}=1\given D_k=0,A]<0.1$。例如，若 $D_k$ 代表肺癌，$k$ 可以年为单位；若 $D_k$ 代表普通感冒病毒感染，$k$ 可以日为单位。时间单位越短，用于拟合 logistic 模型的人时数据集行数就越多。
\end{techpoint}

我们的目标，是估计若没有人在 $k_{\text{end}}$ 之前被删失时本可观测到的生存曲线，其中 $k_{\text{end}}$ 是研究中最长的行政删失时间。也就是说，我们的目标是估计：若即便在删失之后也能知道时变指示变量 $D_k$ 的取值，本可观测到的生存 $\Pr[D_k=0\given A=a]$。严格说来，这一量也可记为 $\Pr[D_k^{\bar c=\bar 0}=0\given A=a]$，其中 $\bar c=(c_1,c_2,\dots,c_{k_{\text{end}}})$。如第 12 章所述，用上标 $\bar c=\bar 0$ 是为了明确我们所设想的量。当不会引起混淆时，我们有时省略上标 $\bar c=\bar 0$。为简单起见，假设每个人的随访开始时间都像是随机指派的——若任何变量都不存在长期趋势，情形便是如此。那么行政删失时间、从而指示变量 $C$，就与治疗和死亡时间都独立。

我们不能仅通过计算接受治疗水平 $a$、存活、且截至 $k$ 未被删失的个体所占比例，来有效估计 $k$ 时的生存 $\Pr[D_k=0\given A=a]$。这个比例是联合概率 $\Pr[C_{k+1}=0,D_{k+1}=0\given A=a]$ 的有效估计量，而后者并非我们想要的量。为看清原因，考虑一项 $k_{\text{end}}=2$ 的研究，其中发生如下情形：
\begin{itemize}
\item $\Pr[C_1=0]=1$，即在 $k=1$ 之前无人被删失；
\item $\Pr[D_1=0\given C_0=0]=0.9$，即 90\% 的个体活过 $k=1$；
\item $\Pr[C_2=0\given D_1=0,C_1=0]=0.5$，即存活者中随机有一半在 $k=2$ 之前被删失；
\item $\Pr[D_2=0\given C_2=0,D_1=0,C_1=0]=0.9$，即其余个体中有 90\% 活过 $k=2$。
\end{itemize}
$k=2$ 时未被删失的存活者比例为 $1\times0.9\times0.5\times0.9=0.405$。然而，若无人被删失，即若 $\Pr[C_2=0\given D_1=0,C_1=0]=1$，则生存本会是 $1\times0.9\times1\times0.9=0.81$。这个例子说明，为什么正确估计生存 $\Pr[D_k=0\given A=a]$ 需要上一节所述的程序。具体而言，在（仿佛）随机指派的删失下，$k$ 时的生存 $\Pr[D_k=0\given A=a]$ 为
\[
\prod_{m=1}^{k}\Pr[D_m=0\given D_{m-1}=0,C_m=0,A=a]\qquad\text{对 }k<k_{\text{end}}
\]
估计程序与上述相同，只是我们要么使用特定原因风险率 $\Pr[D_{k+1}=1\given D_k=0,C_{k+1}=0,A=a]$ 的非参数估计，要么为其拟合 logistic 模型。

我们往往并不准备假设删失是仿佛随机指派的。存在错时进入时，个体的行政删失时间取决于进入研究时的日历时间（较晚进入者的 $k_{\text{end}}$ 较短），而日历时间本身可能与结局相关。因此，上述程序需要校正基线日历时间。此外，还可能有其他基线预后因素在受治疗者（$A=1$）与未受治疗者（$A=0$）之间分布不均，同样需要校正。随后几节把上述程序推广，以通过 g 方法纳入对基线混杂因子的校正。第三部分再把该程序推广到时变治疗与时变混杂因子的情形。

\section{边缘结构模型的逆概率加权}

当受治疗者与未受治疗者不可交换时，直接对比二者的生存曲线无法赋予因果解释。在戒烟例子中，我们估计 120 个月时的生存，戒烟者低于未戒烟者（76.2\% 对 82.0\%），但这未必意味着戒烟会增加死亡。年长者更可能戒烟，也更可能死亡。这种由年龄造成的混杂使戒烟显得有害，因为戒烟者中年长者的比例高于未戒烟者。

定义 $D_k^{a,\bar c=0}$ 为治疗水平 $a$ 且无删失时，$k$ 时死亡的（反事实）时变指示变量。为记号简便，当本章中目标明确是估计无删失情形下的生存时，我们把它写作 $D_k^a$。同样为简便，本章余下部分在 $k$ 时风险率 $\Pr[D_{k+1}=0\given D_k=0,L=l,A]$ 的条件事件中省略 $C_k=0$。也就是说，我们写出的所有表达式都仿佛所有个体有一个共同的行政删失时间，如同我们的戒烟例子。

假设我们要比较若所有人都接受治疗（$a=1$）与都不接受治疗（$a=0$）时本可观测到的反事实生存 $\Pr[D_{k+1}^{a=1}=0]$ 与 $\Pr[D_{k+1}^{a=0}=0]$。也就是说，所关心的因果对比是
\[
\Pr[D_{k+1}^{a=1}=0]\quad\text{vs.}\quad\Pr[D_{k+1}^{a=0}=0]\qquad\text{对 }k=0,2,\dots,k_{\text{end}}-1
\]
由于混杂，前几节所述的生存 $\Pr[D_{k+1}=0\given A=1]$ 与 $\Pr[D_{k+1}=0\given A=0]$ 之对比，未必能有效估计这一对比。要有效估计 $a=1$ 与 $a=0$ 时的量 $\Pr[D_{k+1}^a=0]$，通常需要校正混杂因子，这可以通过若干方法实现。本节聚焦逆概率加权。设受治疗者（$A=1$）与未受治疗者（$A=0$）在 $L$ 各变量取值层内可交换，如图 17.5 的因果图所示。与第 12 至 15 章一样，$L$ 包含性别、年龄、种族、教育、吸烟强度与年限、日常体力活动、休闲锻炼与体重。我们还假定正值性与一致性。逆概率加权生存曲线的估计分两步。\mnote{图 17.5}

\begin{figure}[htbp]
\centering
\begin{tikzpicture}[causal,every node/.style={font=\small}]
  \node (L)  at (-1.9,0)  {$L$};
  \node (A)  at (0,0)     {$A$};
  \node (D1) at (2.2,0)   {$D_1$};
  \node (D2) at (4.4,0)   {$D_2$};
  \node (U)  at (-1.9,-2.0) {$U$};
  \draw[->] (L) -- (A);
  \draw[->] (A) -- (D1);
  \draw[->] (D1) -- (D2);
  \draw[->] (A) to[out=55,in=125] (D2);
  \draw[->] (U) -- (L);
  \draw[->] (U) -- (D1);
  \draw[->] (U) -- (D2);
\end{tikzpicture}
\caption{}
\label{fig:17.5}
\end{figure}

第一步，为研究总体中的每个个体估计稳定化逆概率权重 $SW^A$。程序与第 12 章所述完全相同。我们对给定 $L$ 中各变量时接受治疗（即戒烟）的条件概率 $\Pr[A=1\given L]$ 拟合一个 logistic 模型。所估计的 $SW^A$ 的分母，对受治疗者为 $\widehat{\Pr}[A=1\given L]$，对未受治疗者为 $1-\widehat{\Pr}[A=1\given L]$，其中 $\widehat{\Pr}[A=1\given L]$ 是 logistic 模型的预测值。所估计权重的分子，对受治疗者为 $\widehat{\Pr}[A=1]$，对未受治疗者为 $1-\widehat{\Pr}[A=1]$，其中 $\widehat{\Pr}[A=1]$ 可用非参数方法估计，或取治疗边缘概率 $\Pr[A=1]$ 的 logistic 模型的预测值。预测值的细节见第 11 章。\mnote{代码：程序 17.3}

应用所估计的权重 $SW^A$，会造出一个伪总体，其中 $L$ 中诸变量与治疗 $A$ 独立，从而消除这些变量造成的混杂。在本例中，权重的均值为 1（如所预期），取值范围为 0.33 到 4.21。

第二步，使用人时数据格式，拟合一个与上述类似的风险率模型，只是个体要按其估计的 $SW^A$ 加权。严格说来，这个逆概率加权 logistic 模型估计的是如下边缘结构 logistic 模型的参数
\[
\operatorname{logit}\Pr[D_{k+1}^a=0\given D_k^a=0]=\beta_{0,k}+\beta_1a+\beta_2a\times k+\beta_3a\times k^2
\]
也就是说，逆概率加权模型估计的是：若研究总体中所有人都受治疗（$a=1$）时本可观测到的时变风险率，以及若所有人都不受治疗（$a=0$）时的时变风险率。

随后，可以把由逆概率加权风险率模型得到的 $\Pr[D_{k+1}^a=0\given D_k^a=0]$ 的估计随时间相乘（见上一节），得到治疗 $a=1$ 与无治疗 $a=0$ 下本可观测到的生存 $\Pr[D_{k+1}^a=0]$ 的估计。所得曲线见图 17.6。

\begin{figure}[htbp]
\centering
\begin{tikzpicture}[x=1cm,y=1cm,causal]
  \survaxes
  \draw[thick] (0.00,5.00) -- (0.90,4.90) -- (1.80,4.65) -- (2.70,4.45) -- (3.60,4.25) -- (4.50,4.00) -- (5.40,3.75) -- (6.30,3.58) -- (7.20,3.40) -- (8.10,3.22) -- (9.00,3.07);
  \draw[thick,dashed] (0.00,5.00) -- (0.90,4.86) -- (1.80,4.60) -- (2.70,4.38) -- (3.60,4.10) -- (4.50,3.84) -- (5.40,3.60) -- (6.30,3.45) -- (7.20,3.30) -- (8.10,3.18) -- (9.00,3.05);
  \node[font=\small] at (7.35,4.05) {$A=0$};
  \node[font=\small] at (7.35,2.90) {$A=1$};
\end{tikzpicture}
\caption{}
\label{fig:17.6}
\end{figure}

在本例中，120 个月时的生存估计，在戒烟下为 80.7\%，在未戒烟下为 80.5\%；\mnote{图 17.6}差 0.2\%（基于 500 个自助样本的 95\% 置信区间为 $-4.1\%$ 到 $3.7\%$）。尽管在随访的大部分时间里，治疗下的生存曲线低于无治疗下的曲线，但差值最大也只到 $-1.4\%$（95\% 置信区间 $-3.4\%$ 到 $0.7\%$）。也就是说，通过逆概率加权校正协变量 $L$ 之后，我们在随访期内的任何时点都几乎没有发现戒烟影响死亡率的证据。这一程序的有效性要求治疗模型与边缘风险率模型都无错误设定。

\section{参数 g 公式}

上一节中，我们通过逆概率加权估计了整个研究总体上治疗下与无治疗下的生存曲线。为此，我们校正了 $L$，并假定了可交换性、正值性与一致性。在这些假定下估计边缘生存曲线的另一种方法，是基于参数模型的标准化，即参数 g 公式。

治疗水平 $a$ 下 $k+1$ 时的生存 $\Pr[D_{k+1}^a=0]$，是协变量 $L$ 与治疗水平 $A=a$ 各层内 $k+1$ 时条件生存概率的加权平均，权重为 $L$ 各水平 $l$ 上的个体比例。也就是说，在可交换性、正值性与一致性下，$\Pr[D_{k+1}^a=0]$ 等于标准化生存
\[
\sum_l\Pr[D_{k+1}=0\given L=l,A=a]\Pr[L=l]
\]
形式化的证明见 2.3 节。

因此，参数 g 公式的估计分两步。第一步，用我们的行政删失数据估计条件生存 $\Pr[D_{k+1}=0\given L=l,A=a]$。第二步，对协变量 $L$ 的所有取值 $l$ 计算其加权平均。下面结合戒烟例子说明这两步。

第一步，拟合一个与 17.2 节所述类似的参数风险率模型，只是把 $L$ 中各变量作为协变量纳入。若模型设定正确，它就能有效估计治疗 $A$ 与协变量 $L$ 各层内的时变风险率 $\Pr[D_{k+1}=1\given D_k=0,L,A]$。“1 减条件风险率”之积
\[
\prod_{m=0}^{k}\Pr[D_{m+1}=0\given D_m=0,L=l,A=a]
\]
等于条件生存 $\Pr[D_{k+1}=0\given L=l,A=a]$。由于给定 $L$ 的条件可交换性，特定协变量取值 $L=l$ 与 $A=a$ 下的条件生存，可以因果地解释为：若具备该组协变量的所有人都接受治疗取值 $a$ 时本可观测到的生存。也就是说，
\[
\Pr[D_{k+1}=0\given L=l,A=a]=\Pr[D_{k+1}^a=0\given L=l]
\]
因此，条件风险率可用于估计 $L$ 的每一组取值 $l$ 下、治疗（$a=1$）与无治疗（$a=0$）时的生存曲线。例如，我们可以用这个模型估计 61 岁、受过大学教育、锻炼水平低等等的白人男性的治疗与无治疗生存曲线。然而，我们的目标是估计边缘的、而非条件的治疗与无治疗生存曲线。\mnote{在第 12 章中，我们把对全部协变量 $L$ 取条件的模型称为伪边缘结构模型。}

第二步，对协变量 $L$ 的所有取值 $l$ 计算条件生存的加权平均，即把生存标准化到混杂因子分布上。为此，我们使用 13.3 节所述的均值标准化方法：先扩展数据集，再对结局建模并预测，然后取平均。即便 $L$ 中某些变量是连续的（此时对 $l$ 求和形式上是一个积分），该方法仍然可用。所得曲线见图 17.7。\mnote{代码：程序 17.4。该程序与第 13 章所述的程序类似。}

\begin{figure}[htbp]
\centering
\begin{tikzpicture}[x=1cm,y=1cm,causal]
  \survaxes
  \draw[thick] (0.00,5.00) -- (0.90,4.90) -- (1.80,4.66) -- (2.70,4.46) -- (3.60,4.25) -- (4.50,4.02) -- (5.40,3.78) -- (6.30,3.58) -- (7.20,3.38) -- (8.10,3.20) -- (9.00,3.06);
  \draw[thick,dashed] (0.00,5.00) -- (0.90,4.85) -- (1.80,4.60) -- (2.70,4.37) -- (3.60,4.08) -- (4.50,3.84) -- (5.40,3.62) -- (6.30,3.45) -- (7.20,3.28) -- (8.10,3.15) -- (9.00,3.04);
  \node[font=\small] at (7.35,4.05) {$A=0$};
  \node[font=\small] at (7.35,2.90) {$A=1$};
\end{tikzpicture}
\caption{}
\label{fig:17.7}
\end{figure}

在本例中，治疗下的生存曲线在整个随访期内都低于无治疗下的曲线，\mnote{图 17.7}但差值最大也只到 $-2.0\%$（95\% 置信区间 $-5.6\%$ 到 $1.8\%$）。120 个月时的生存估计，在戒烟下为 80.4\%，在未戒烟下为 80.6\%；差 0.2\%（95\% 置信区间 $-4.6\%$ 到 $4.1\%$）。也就是说，通过标准化校正协变量 $L$ 之后，我们在随访期内的任何时点都几乎没有发现戒烟影响死亡率的证据。注意，通过逆概率加权（上一节）与参数 g 公式（本节）估计的生存曲线相似但不完全相同，因为二者依赖不同的参数假定：逆概率加权估计要求治疗模型与无条件风险率模型都无错误设定；参数 g 公式估计则要求条件风险率模型无错误设定。

\section{结构嵌套模型的 g 估计}

前几节描述的因果对比，比较的是不同治疗水平 $A$ 下的生存或风险。生存是由 logistic 回归模型估计的风险率算出的。当分析方法是边缘结构模型的逆概率加权或参数 g 公式时，这一做法可行；但当方法是结构嵌套模型的 g 估计时则不可行。如第 14 章所解释，结构嵌套模型是关于条件因果对比（如不同治疗水平下协变量特定均值之差或之比）的模型，而不是关于这些对比的各个组成部分（如不同治疗水平下的各个均值）的模型。因此，我们无法用结构嵌套模型估计生存或风险率。\mnote{事实上，我们甚至可能无法近似风险率比，因为结构嵌套 logistic 模型不易推广到时变治疗（技术注记 14.1）。}

不过，我们可以考虑用结构嵌套对数线性模型来建模不同治疗水平下累积发病（即风险）之比。结构嵌套累积失效时间模型正是这么做的（见技术注记 17.2），但当失效是罕见事件时，用它们最为合适，因为对数线性模型不会自然地给风险设定 1 这一上限。对非罕见的失效，我们可以转而考虑用结构嵌套对数线性模型来建模不同治疗水平下累积生存概率（即 1 减去风险）之比。结构嵌套累积生存时间模型正是这么做的（见技术注记 17.2），但当生存是罕见事件时，用它们最为合适，因为对数线性模型不会自然地给生存设定 1 这一上限。更一般的选择，是用结构嵌套模型来建模不同治疗方案下生存时间之比。也就是说，加速失效时间（AFT）模型。\mnote{Tchetgen Tchetgen 等（2015）与 Robins（1997b）描述了用工具变量做生存分析，其问题与这里针对结构嵌套模型所述的问题类似。}

设 $T_i^a$ 为个体 $i$ 在治疗水平 $a$ 下的反事实生存时间。处理 $A$ 对个体 $i$ 生存时间的影响，可用其治疗下与无治疗下反事实生存时间之比 $T_i^{a=1}/T_i^{a=0}$ 来度量。若生存时间比大于 1，则治疗有益，因为它延长了生存时间；若比值小于 1，则治疗有害；若比值等于 1，则治疗无效应。暂且假设处理对总体中每个人的效应都相同。

于是我们可以考虑一个结构嵌套加速失效时间（AFT）模型，其形式为 $T_i^a/T_i^{a=0}=\exp(-\psi_1a)$，其中 $\psi_1$ 度量处理所致每个个体生存时间的扩张（或收缩）。若 $\psi_1<0$ 则治疗延长生存时间，若 $\psi_1>0$ 则治疗缩短生存时间，若 $\psi_1=0$ 则治疗不影响生存时间。\mnote{$\psi$ 前的负号保留了通常的解释：正参数表示有害，负参数表示有益。}

\begin{techpoint}[结构嵌套累积失效时间（CFT）模型与累积生存时间（CST）模型]
对时不变处理，一个（非嵌套的）结构 CFT 模型是这样一种模型：给定处理 $A$ 与协变量 $L$，它刻画治疗取值 $a$ 下的反事实风险与治疗取值 0 下的反事实风险之比。模型的一般形式为
\[
\frac{\Pr[D_k^a=1\given L,A]}{\Pr[D_k^{a=0}=1\given L,A]}=\exp[\gamma_k(L,A;\psi)]
\]
其中 $\gamma_k(L,A;\psi)$ 是处理与协变量历史的函数，由（可能是向量值的）参数 $\psi$ 索引。为满足一致性，当 $A=0$ 时 $\exp[\gamma_k(L,A;\psi)]$ 必须为 1，因为此时所比较的两个治疗取值相同；当 $m$ 时的处理对 $k$ 时的结局没有效应时，它也必须为 1。这类函数的一个例子是 $\gamma_k(L,A;\psi)=\psi A$，于是 $\psi=0$ 对应无效应，$\psi<0$ 对应有益效应，$\psi>0$ 对应有害效应。

类似地，对时不变处理，一个（非嵌套的）结构 CST 模型是这样一种模型：给定处理 $A$ 与协变量 $L$，它刻画治疗取值 $a$ 下的反事实生存与治疗水平 0 下的反事实生存之比。模型的一般形式为
\[
\frac{\Pr[D_k^a=0\given L,A]}{\Pr[D_k^{a=0}=0\given L,A]}=\exp[\gamma_k(L,A;\psi)]
\]
尽管 CFT 与 CST 模型的区别只在于我们是对 $\Pr[D_k^a=1\given L,A]$ 还是对 $\Pr[D_k^a=0\given L,A]$ 设定乘法模型，但 $\gamma_k(L,A;\psi)$ 的含义不同，因为只要处理效应非零，关于风险的乘法模型就不是关于生存的乘法模型。当我们让时间下标 $k$ 连续而非离散时，结构 CST 模型等价于结构加法风险率模型（Tchetgen Tchetgen 等，2015）：任何关于 $\Pr[D_k^a=0\given L,A]/\Pr[D_k^{a=0}=0\given L,A]$ 的模型，都会诱导出关于各时点特定风险率之差的一个模型，反之亦然。

使用结构 CFT 模型要求：对协变量 $L$ 的所有取值，全部治疗取值下的条件累积失效概率都满足某类罕见失效假定。在这种“罕见失效”情形下，结构 CFT 模型比 AFT 模型有一项优势：它容许无偏的估计方程，这些方程对模型参数可微，因而易于求解。Page（2005）与 Picciotto 等（2012）对结构 CFT 与 CST 模型给出了更多细节。对时变处理，这一类模型可视为多元结构嵌套均值模型（Robins 1994）的特例。见技术注记 14.1。
\end{techpoint}

更一般地，处理效应可能依赖于协变量 $L$，因此更一般的结构 AFT 模型为 $T_i^a/T_i^{a=0}=\exp(-\psi_1a-\psi_2aL_i)$，其中 $\psi_1$ 与 $\psi_2$（一个向量）在个体间为常数。重新整理各项，模型可写为
\[
T_i^{a=0}=T_i^a\exp(\psi_1a+\psi_2aL_i)\qquad\text{对全部个体 }i
\]

沿用第 14 章的推理，反事实的一致性蕴含模型 $T_i^{a=0}=T_i\exp(\psi_1A_i+\psi_2A_iL_i)$，其中反事实时间 $T_i^a$ 被实际生存时间 $T_i^{A_i}=T_i$ 替换。参数 $\psi_1$ 与 $\psi_2$ 可用本节后文描述的、经修改的 g 估计程序（以处理行政删失）来估计。

上述结构 AFT 模型并不现实，因为它既是确定性的，又是保秩的。说它确定性，是因为它假定对每个个体，无治疗下的反事实生存时间 $T^{a=0}$ 都可以无误差地由观测生存时间 $T$、治疗 $A$ 与协变量 $L$ 的函数算出。说它保秩，是因为在该模型下，若个体 $i$ 在两人都未治疗时会先于个体 $j$ 死亡，即 $T_i^{a=0}<T_j^{a=0}$，则个体 $i$ 在两人都受治疗时也会先于个体 $j$ 死亡，即 $T_i^{a=1}<T_j^{a=1}$。

由于保秩假设不可信，一般不应使用依赖它的因果推断方法，这一点我们在第 14 章讨论过。然而我们在此仍将使用保秩模型来描述结构 AFT 模型的 g 估计，因为对保秩模型而言 g 估计更易理解，而且保秩与非保秩模型的 g 估计程序实际上相同。\mnote{Robins（1997b）描述了非确定性、非保秩的结构嵌套 AFT 模型。}

考虑不含处理与协变量乘积项的较简单保秩模型 $T_i^{a=0}=T_i\exp(\psi A_i)$。若不存在行政删失，即若我们能观测到所有人的死亡时间 $T$，则该结构 AFT 模型参数 $\psi$ 的 g 估计将十分直接。事实上，那种情形下的 g 估计程序与 14.5 节所述相同。第一步，在许多可能的因果参数取值 $\psi^\dagger$ 下计算候选反事实 $H_i(\psi^\dagger)=T_i\exp(\psi^\dagger A_i)$。\mnote{计算强度较低的方法称为定向搜索法，统计软件中提供了用于近似搜索的实现。Nelder-Mead 单纯形法就是一种定向搜索法。}第二步，找出这样的 $\psi^\dagger$：使所得的 $H_i(\psi^\dagger)$ 在以 $H_i(\psi^\dagger)$ 与混杂因子 $L$ 为协变量、对 $A=1$ 的概率拟合的 logistic 模型中与治疗 $A$ 独立。这样的 $\psi^\dagger$ 就是 $\psi$ 的 g 估计值。

然而，在存在时间 $K$ 处行政删失的情况下，这一程序无法实施，因为对 $T_i$ 未知的个体无法计算 $H_i(\psi^\dagger)$。人们或许会想把 g 估计程序限制在生存时间可观测的个体上，即 $T_i\le K$ 者。不幸的是，这种做法会导致选择偏倚。为看清原因，考虑下面这个过度简化的情景。

我们做一项为期 60 个月的随机化实验，估计二值处理 $A$ 对生存时间 $T$ 的效应。研究只有 3 类个体参与。1 类个体是那些在无治疗下会在 36 个月时死亡的人（$T^{a=0}=36$）。2 类个体是那些在无治疗下会在 72 个月时死亡的人（$T^{a=0}=72$）。3 类个体是那些在无治疗下会在 108 个月时死亡的人（$T^{a=0}=108$）。也就是说，3 类个体预后最好，1 类个体预后最差。由于随机化，我们预期 1 类、2 类与 3 类个体在两个治疗组 $A=1$ 与 $A=0$ 中的比例相同。也就是说，受治疗者与未受治疗者预期可交换。

\begin{table}[htbp]
\centering
\caption{}
\label{tab:17.1}
\begin{tabular}{lccc}
\toprule
类型 & 1 & 2 & 3\\
\midrule
$T^{a=0}$ & 36 & 72 & 108\\
$T^{a=1}$ & 24 & 48 & 72\\
\bottomrule
\end{tabular}
\end{table}

设处理 $A=1$ 相对于 $A=0$ 缩短生存时间。\mnote{表 17.1}表 17.1 给出了每类个体在治疗下与无治疗下的生存时间。由于随访的行政终止为 $K=60$ 个月，1 类个体的死亡无论被随机分到 $A=1$ 还是 $A=0$ 都会被观测到（两个生存时间都小于 60），3 类个体的死亡无论被分到 $A=1$ 还是 $A=0$ 都会被行政删失（两个生存时间都大于 60）。然而，2 类个体的死亡只有在被分到 $A=1$ 时才会被观测到。因此，一项纳入所有死亡时间未被行政删失者的分析，会在受治疗者与未受治疗者之间造成个体类型的不平衡。可交换性将被破坏，因为 $A=1$ 组会包含 1 类与 2 类个体，而 $A=0$ 组只包含 1 类个体。$A=0$ 组的个体平均预后比 $A=1$ 组差，这会使治疗显得比实际更好。当处理对生存时间有非零效应时，这种选择偏倚（第 8 章）就会出现。

\begin{techpoint}[人为删失]
设 $K(\psi)$ 为可能与一个实际上在时间 $K$（随访的行政终止）死亡的个体相对应的、无治疗下的最小生存时间。对二值处理 $A$，$K(\psi)=\inf\{K\exp(\psi A)\}$，这意味着：若治疗缩短生存时间（即 $\psi>0$），则 $K(\psi)=K\exp(\psi\times0)=K$；若治疗延长生存时间（即 $\psi<0$），则 $K(\psi)=K\exp(\psi\times1)=K\exp(\psi)$；若治疗不影响生存时间（即 $\psi=0$），则 $K(\psi)=K\exp(0)=K$。

所有被行政删失者（即 $T>K$）都有 $\Delta(\psi)=0$，因为至少存在一个治疗水平（即他们实际接受的那个）使其生存时间大于 $K$，即 $H(\psi)\ge K(\psi)$。一些未被行政删失者（即 $T\le K$）也有 $\Delta(\psi)=0$，他们被排除在分析之外——即被人工删失——以避免选择偏倚。

人为删失指示变量 $\Delta(\psi)$ 是 $H(\psi)$ 与 $K$ 的函数。在给定 $L$ 的条件可交换性下，所有这类函数在 $\psi$ 取真值时，都给定协变量 $L$ 而与治疗 $A$ 条件独立。也就是说，AFT 模型参数的 g 估计可以基于 $\Delta(\psi)$ 而非 $H(\psi)$ 来进行。严格说来，是用 $\Delta(\psi)$ 替换技术注记 14.2 之估计方程中的 $H(\psi)$。实际估计细节见 Hernán 等（2005）的附录。
\end{techpoint}

为避免这种选择偏倚，需要选取这样的个体：无论他们受治疗还是未受治疗，其生存时间本都会在随访结束时被观测到，即 $T_i^{a=0}\le K$ 且 $T_i^{a=1}\le K$ 者。在本例中，为保持可交换性，我们将不得不把所有 2 类个体排除在分析之外。也就是说，我们不仅排除 $T_i>K$ 的行政删失个体，还要排除一些 $T_i\le K$、生存时间已知的未删失个体，因为若他们接受的不是实际所接受的治疗水平，其生存时间本会大于 $K$。\mnote{这种把未删失个体排除在分析之外的做法，常称为人为删失。见技术注记 17.3。}

随后我们定义一个指示变量 $\Delta(\psi)$：个体被排除时取 0，不被排除时取 1。于是 g 估计程序经修改，把变量 $H(\psi^\dagger)$ 换成指示变量 $\Delta(\psi^\dagger)$。细节见技术注记 17.3。在本例中，保秩结构 AFT 模型 $T_i^{a=0}=T_i\exp(\psi A_i)$ 的 g 估计值 $\hat\psi$ 为 $-0.047$（95\% 置信区间：$-0.223$ 到 $0.333$）。$\exp(-\hat\psi)=1.05$ 可解释为：若研究中所有人都接受 $a=1$ 时本可观测到的中位生存时间，除以若所有人都接受 $a=0$ 时本可观测到的中位生存时间。这一生存时间比提示戒烟 $A$ 对死亡时间几乎没有影响。\mnote{代码：程序 17.5。$\psi$ 的点估计是使技术注记 17.3 所述估计函数取最小值的那个值；95\% 置信区间的界限是使该估计函数取 3.84（自由度为 1 的 $\chi^2$）的那些值。}

正如我们在第 14 章所说，结构嵌套模型（包括 AFT 模型）在实践中很少被使用。实施该方法的一个实际障碍是缺少易用的软件。在生存分析情形下，一个更严重的障碍是：结构 AFT 模型的参数需要通过搜索算法估计，而搜索算法并不保证找到唯一解。对含两个或更多参数 $\psi$ 的模型，这个问题更加严重。因此，该方法已发表的少数应用往往使用不允许处理效应随协变量取值变化的简单 AFT 模型。

这种状况令人遗憾，因为专业知识（如生物学机制）更容易转译为结构 AFT 模型的参数，而不是结构风险率模型的参数。在使用非确定性、非保秩的结构 AFT 模型时，尤其如此。
